\documentclass[11pt]{article}

\usepackage[margin=1in]{geometry}
\usepackage{graphicx}
\usepackage{booktabs}
\usepackage{amsmath}
\usepackage{hyperref}

\title{CanopyGuard-AI: Vegetation Risk Forecasting and Maintenance Prioritization Near Power-Line Corridors}
\author{CanopyGuard-AI contributors}
\date{}

\begin{document}

\maketitle

\begin{abstract}
Vegetation encroachment near power-line corridors creates maintenance and
reliability challenges. This paper presents CanopyGuard-AI, a reproducible
pipeline for forecasting vegetation risk and prioritizing trimming decisions
using satellite time series, sparse canopy reference data, LiDAR validation, and
optimization. The risk score is designed as a research prioritization aid and is
not a regulatory clearance or safety certification.
\end{abstract}

\section{Introduction}

Power-line vegetation management requires identifying where growth is likely to
create future risk and deciding how limited maintenance resources should be
allocated. This work studies whether remote sensing and forecasting can produce
a defensible prioritization signal for vegetation management.

\section{Related Work}

This section will summarize prior work on vegetation management, satellite-based
canopy monitoring, GEDI and LiDAR validation, forecasting, and maintenance
optimization.

\section{Study Area and Data}

This section will describe the study area, HIFLD corridor context, Sentinel-2
time series, GEDI reference data, and LiDAR validation data.

\section{Methods}

This section will define the feature pipeline, forecasting task, risk score, and
scheduling optimizer.

\section{Experiments}

This section will define time-based train, validation, and test splits,
baselines, metrics, and ablation studies.

\section{Results}

This section will report forecasting accuracy, risk-prioritization performance,
scheduling outcomes, and uncertainty analysis.

\section{Discussion}

This section will interpret results, failure modes, operational constraints, and
scientific limitations.

\section{Conclusion}

This section will summarize the contribution and identify next validation steps.

\section*{AI Assistance Disclosure}

AI assistance was used for project planning, code drafting, and manuscript
support. The final manuscript, claims, results, and citations must be verified by
the authors according to the target venue policy.

\end{document}

